%
% hub_display_reconciliation.tex
%
% A Z specification of the DES-068 reconciliation policy: single-owner
% preemption of the Hub's declared identity on the Hub-Display socket leg,
% and the manifest-driven purge that follows every fresh identify.
%
% Models the DECISION in
% docs/architecture/hub-display-reconciliation-design.md, section "Decision:
% manifest-on-identify with single-owner preemption", as implemented by
% HubReconciliation (display/hub_reconciliation.py) and SceneReplica
% .scenes_to_purge (display/replica/scene_replica.py).
%
% Type-check:  fuzz docs/hub_display_reconciliation.tex
% Model-check (see the Verification section for the exact goal predicates):
%   probcli docs/hub_display_reconciliation.tex -model_check \
%       -p DEFAULT_SETSIZE 3
%
% Formatting follows the fuzz house rules: \quad~ for continuation indent,
% schema lines under 80 columns, predicates broken at \land/\lor/\implies.
%

\documentclass[a4paper,10pt,fleqn]{article}
\usepackage[margin=1in]{geometry}
\usepackage{fuzz}
\usepackage[colorlinks=true,linkcolor=blue,citecolor=blue,urlcolor=blue]{hyperref}

\begin{document}

\title{Hub/Display Scene Reconciliation: a Z Specification}
\author{Formal model of DES-068's single-owner preemption and manifest purge}
\date{August 2026}
\maketitle

% ============================================================
\section{Introduction}
% ============================================================

Killing and restarting \texttt{luxd} drops its in-memory scene store
entirely; the Display was never told to forget anything, so it keeps
rendering the dead process's scenes forever --- a click on one dispatches
into an owner file descriptor that maps to nothing and silently does
nothing.  DES-068 closes this with two mechanisms working together: a
\texttt{kind="hub"} connection's identify forcibly evicts any predecessor
declaring the same identity (\emph{single-owner preemption}), and every
such connection sends a manifest of its complete live-scene set
immediately after identifying, which the Display uses to purge every
scene it cannot attribute to that manifest (\emph{manifest-driven purge}).

The two mechanisms exist for the same reason: a manifest alone does not
close the interleaving where a straggling \texttt{SceneMessage}, still
buffered in the OS receive queue from the \emph{old}, about-to-be-superseded
connection, is processed \emph{after} the new connection's manifest has
already purged the scene it names --- silently re-materializing exactly the
ghost the manifest just told the Display to forget.  Preemption closes
that hole at the socket layer: \texttt{SocketListener.remove\_client}
closes the predecessor's socket and drops it from the poll set, so
\texttt{poll\_clients} can never again read a frame from it, and whatever
bytes were still buffered on it are discarded by the OS.  That is a fact
about which \emph{sockets are still open}, not about which connection is
currently declared the Hub's identity --- the two are usually the same
connection, but the whole hazard this design closes is the interleaving
where they briefly are not.  This document therefore tracks both:
\texttt{hubOwner}, the declared identity a manifest is trusted against, and
\texttt{liveConns}, which sockets can still deliver a message at all.

Four properties must hold across every interleaving:

\begin{description}
  \item[I1 --- no stale owner.] Every scene the Display holds is either
    attributed to the live \texttt{hubOwner}, or is orphaned and still
    named in the most recent manifest (legitimately awaiting reclaim, not
    a ghost).
  \item[I2 --- only-live-hub-installs.] A scene install is accepted only
    from a connection whose socket is still open.  An install from a
    connection preemption has already evicted is dropped, not applied ---
    this is the ordering property single-owner preemption exists to close.
  \item[I3 --- no two hub winners.] At every state, at most one connection
    holds the Hub's declared identity.
  \item[I4 --- deadlock freedom.] No interleaving traps the model in a
    state where an old and new claimant are both partially processed.
\end{description}

% ============================================================
\section{Basic Types}
% ============================================================

\begin{zed}
  [CONNECTION\_ID, SCENE\_ID]
\end{zed}

\texttt{CONNECTION\_ID} is the set of Hub-Display socket connections the
Display can hold open; three suffice to model the old Hub fd, the new Hub
fd, and one unrelated \texttt{kind="test"} connection that must be
unaffected by preemption.  \texttt{SCENE\_ID} is the set of scenes the
Display can hold; two suffice to exhibit a mixed frame where one scene
survives a manifest and its sibling does not.  Both carriers are bounded
by ProB's \texttt{DEFAULT\_SETSIZE}.

% ============================================================
\section{Free Types}
% ============================================================

\begin{zed}
  OWNER ::= orphan | conn \ldata CONNECTION\_ID \rdata
\end{zed}

The attribution of one scene on the Display.  \texttt{orphan} is the
sentinel \texttt{scene\_to\_owner} carries after its owning connection
disconnects (\texttt{reassign\_scenes\_of}'s real behaviour: reassigned,
never removed).  \texttt{conn(c)} is a scene currently attributed to a
connection \texttt{c}, live or not --- \texttt{scene\_to\_owner} itself
carries no liveness check, which is exactly the gap I2 exists to close one
layer up, at the point a \emph{new} install is accepted.

\begin{zed}
  FLAG ::= set | clear
\end{zed}

A two-valued marker, used for \texttt{manifestDeclared} below.  A
dedicated free type is used rather than a built-in boolean so the spec is
self-contained for \texttt{fuzz}, matching \texttt{hub\_replicator.tex}'s
own \texttt{FLAG}.

% ============================================================
\section{State}
% ============================================================

The state is flat: one schema.  Neither \texttt{hubOwner} nor
\texttt{liveConns} carries a cardinality bound here --- ``at most one Hub
owner'' is I3, and ``a preempted connection cannot still deliver'' is I2,
the very properties single-owner preemption exists to establish.  Bounding
either in the schema would enforce it as a guard on every successor,
silently disabling any operation that could break it and masking the
defect the fidelity variant (Section~\ref{sec:fidelity}) must be able to
reach.

\begin{schema}{Recon}
  hubOwner : \power CONNECTION\_ID \\
  liveConns : \power CONNECTION\_ID \\
  sceneOwner : SCENE\_ID \pfun OWNER \\
  manifested : \power SCENE\_ID \\
  manifestDeclared : FLAG
\end{schema}

\texttt{hubOwner} is the connection currently declaring the Hub's
identity --- empty or a single connection in every correct trace, and the
set ProB may grow past one only if the model itself is broken (no
operation below can grow it past one; I3's check confirms this holds).
\texttt{liveConns} is which connections' sockets are still open enough to
deliver a message --- the \texttt{SocketListener.\_clients} membership
test \texttt{poll\_clients} runs before it ever reaches
\texttt{\_handle\_message}.  \texttt{sceneOwner} is the Display's
\texttt{scene\_to\_owner} map: a partial function whose domain is exactly
the scenes currently on the Display, so purging a scene removes it from
the domain outright, matching \texttt{dismiss\_framed\_scene} discarding
the scene from its frame.  \texttt{manifested} is the most recent
\texttt{HubManifestMessage}'s scene-id set, kept across identifies so the
invariants can still see what the current \texttt{hubOwner} last declared.
\texttt{manifestDeclared} scopes I1 the same way \texttt{hub\_replicator
.tex}'s \texttt{Quiescent} scopes its I4: a freshly identified owner has
not purged anything yet, and a connection dropping can mint a fresh orphan
the last manifest never promised to cover, so either is a legitimate,
transient ``awaiting reclaim'' window, not yet a ghost ---
\texttt{HubIdentify} and \texttt{ConnectionDrop} clear the flag,
\texttt{HubManifest} sets it, \texttt{SceneInstall} leaves it alone (an
owned scene always satisfies I1's first disjunct once its installer is
also the live \texttt{hubOwner}), and I1 is required only once the flag is
set.

% ============================================================
\section{Initialization}
% ============================================================

\begin{schema}{Init}
  Recon'
\where
  hubOwner' = \emptyset \\
  liveConns' = \emptyset \\
  sceneOwner' = \emptyset \\
  manifested' = \emptyset \\
  manifestDeclared' = clear
\end{schema}

A single initial state: no connection holds the Hub's identity or an open
socket, the Display holds no scenes, and no manifest has ever been
declared --- the cold-start condition before any \texttt{luxd} has ever
connected.

% ============================================================
\section{Operations}
% ============================================================

\subsection{HubIdentify --- single-owner preemption}

A connection \texttt{c?} declares \texttt{kind="hub"}.  Any connection
previously in \texttt{hubOwner} is evicted from \emph{both}
\texttt{hubOwner} and \texttt{liveConns} --- \texttt{SocketListener
.remove\_client} closing the predecessor's socket and
\texttt{RenderLoop.\_on\_client\_disconnected} reassigning its scenes in
one call --- and every scene it owned is orphaned, never removed
(\texttt{reassign\_scenes\_of}'s reassign-not-remove behaviour).  \texttt{c?}
itself joins \texttt{liveConns} here too: the model treats ``connect'' and
``identify'' as one atomic step, since no interleaving between a bare TCP
accept and the \texttt{ConnectMessage} that follows it matters to this
design.  \texttt{manifested} is untouched: a fresh identify carries no
manifest of its own until the following \texttt{HubManifest} declares
one, mirroring the wire order (\texttt{ConnectMessage} then
\texttt{HubManifestMessage}, never combined).

\begin{schema}{HubIdentify}
  \Delta Recon \\
  c? : CONNECTION\_ID
\where
  liveConns' = (liveConns \setminus hubOwner) \cup \{ c? \} \\
  sceneOwner' = sceneOwner \oplus \\
  \quad~ (\{ s : \dom sceneOwner | \\
  \quad~ \quad~ (\exists cid : hubOwner @ sceneOwner~s = conn(cid)) \} \cross \{ orphan \}) \\
  hubOwner' = \{ c? \} \\
  manifested' = manifested \\
  manifestDeclared' = clear
\end{schema}

\subsection{HubManifest --- the purge}

Only the current \texttt{hubOwner} may declare a manifest.  The purge is
domain restriction: every scene kept is either named in the new manifest
or still owned by \texttt{c?} itself (rule item 4 of the design ---
ownership self-corrects on the Hub's own next resend, so a scene the
manifest names but \texttt{c?} does not yet own is left exactly as it
was, orphan or otherwise).  Every other scene --- including an orphan the
manifest does not reclaim --- leaves the domain outright.

\begin{schema}{HubManifest}
  \Delta Recon \\
  c? : CONNECTION\_ID \\
  m? : \power SCENE\_ID
\where
  c? \in hubOwner \\
  manifested' = m? \\
  sceneOwner' = \{ s : \dom sceneOwner | \\
  \quad~ s \in m? \lor sceneOwner~s = conn(c?) \} \dres sceneOwner \\
  hubOwner' = hubOwner \\
  liveConns' = liveConns \\
  manifestDeclared' = set
\end{schema}

\subsection{SceneInstall --- only from a live connection}

A scene push is attributed to \texttt{c?} and accepted only when its
socket is still open.  This guard \emph{is} I2: \texttt{HubIdentify}
removes a preempted predecessor from \texttt{liveConns} in the same
atomic step that installs the new owner, so a straggling install already
buffered on the old socket can never fire again once that step has run
--- exactly matching \texttt{remove\_client} closing the fd before
\texttt{poll\_clients} could ever read another frame from it.  This
guard is deliberately \emph{not} \texttt{c? \in hubOwner}: the real
\texttt{\_handle\_scene} performs no ownership check at all, Hub-kind or
otherwise --- any connection whose socket is still open may push a
scene, matching \texttt{scene\_to\_owner} being keyed by raw fd with no
kind distinction.

\begin{schema}{SceneInstall}
  \Delta Recon \\
  c? : CONNECTION\_ID \\
  s? : SCENE\_ID
\where
  c? \in liveConns \\
  sceneOwner' = sceneOwner \oplus \{ s? \mapsto conn(c?) \} \\
  hubOwner' = hubOwner \\
  liveConns' = liveConns \\
  manifested' = manifested \\
  manifestDeclared' = manifestDeclared
\end{schema}

\subsection{ConnectionDrop --- ordinary disconnect}

Models an unprompted disconnect: the old \texttt{luxd} process, or any
other connection, dies on its own, closing its socket before any new
connection ever identifies --- the ordinary case of a restart, where by
the time anyone reconnects the old connection is already gone and
\texttt{HubIdentify}'s own eviction step (Section~\ref{sec:fidelity}
traces the case where a competing identify races it instead) has nothing
left to preempt.  \texttt{c?} leaves both \texttt{hubOwner} (if it was
there) and \texttt{liveConns}, and every scene it owned is orphaned,
never removed --- the same reassign-not-remove rule
\texttt{HubIdentify}'s preemption step uses.

\begin{schema}{ConnectionDrop}
  \Delta Recon \\
  c? : CONNECTION\_ID
\where
  hubOwner' = hubOwner \setminus \{ c? \} \\
  liveConns' = liveConns \setminus \{ c? \} \\
  sceneOwner' = sceneOwner \oplus \\
  \quad~ (\{ s : \dom sceneOwner | sceneOwner~s = conn(c?) \} \cross \{ orphan \}) \\
  manifested' = manifested \\
  manifestDeclared' = clear
\end{schema}

% ============================================================
\section{Invariants}
\label{sec:inv}
% ============================================================

Four properties must hold across all reachable states.  None is placed in
the \texttt{Recon} schema: a predicate placed there would be enforced as a
guard on every successor, silently disabling any operation that would
break it and thereby masking the very defect this model exists to catch.

\paragraph{Invariant 1 --- no stale owner.}
Once the current owner has declared a manifest, and nothing has happened
since to invalidate that declaration, every scene the Display holds is
either owned by the live \texttt{hubOwner}, or is orphaned and still named
in that manifest:
\[
  manifestDeclared = set \implies \\
  \quad~ \forall s : \dom sceneOwner @ \\
  \quad~ \quad~ (\exists cid : hubOwner @ sceneOwner~s = conn(cid)) \\
  \quad~ \quad~ \lor (sceneOwner~s = orphan \land s \in manifested).
\]
\texttt{HubManifest}'s domain restriction is exactly this predicate's
positive case, applied at every manifest: a scene surviving the purge is
either owned by \texttt{c?} (the first disjunct, since
\texttt{c? \in hubOwner} at that point) or named in \texttt{m?} (the
second, once \texttt{manifested' = m?}) --- and \texttt{HubManifest} is the
only operation that sets \texttt{manifestDeclared}.  Unlike a model whose
\texttt{SceneInstall} is gated on \texttt{hubOwner} membership, gating it
on \texttt{liveConns} instead means a genuinely stale install --- one from
a connection \texttt{HubIdentify} has already evicted from
\texttt{liveConns} but that a buggy variant left in \texttt{hubOwner} or
never evicted at all --- surfaces exactly here: its scene is owned by a
connection outside the live \texttt{hubOwner} and is not an orphan, so
neither disjunct holds.  \texttt{HubIdentify} and \texttt{ConnectionDrop}
are the only operations that can introduce an orphan the standing
manifest never promised to cover, and both clear \texttt{manifestDeclared}
when they run, so the guard is never true at a state their own action
just destabilized.

\paragraph{Invariant 2 --- only-live-hub-installs.}
Stated directly as \texttt{SceneInstall}'s guard: no successor state can
be reached by installing a scene under a connection whose socket has
already been closed.  \texttt{HubIdentify}'s preemption step is the only
place a connection's socket closes as a \emph{side effect} of another
connection's action (\texttt{ConnectionDrop} closes only its own
argument's socket), and it does so atomically with recording the new
\texttt{hubOwner} --- there is no intermediate state in which the old
connection is preempted from \texttt{hubOwner} but still live in
\texttt{liveConns}, or the reverse.  A direct reachability check on
``\texttt{sceneOwner} names a connection outside \texttt{liveConns}'' is
structurally tautological in \emph{both} the correct spec and the buggy
variant of Section~\ref{sec:fidelity}: \texttt{ConnectionDrop} orphans
exactly the scenes it drops liveness for, in the same atomic step, so
\texttt{sceneOwner}'s range can never point at a departed connection in
either model --- that relationship was never what the bug breaks.  What
I2's guard actually protects against, and what does distinguish the two
models, is a stale \emph{install still landing} once its source was
supposed to have been evicted; that failure surfaces as an I1 violation
(a scene owned by a connection \texttt{hubOwner} no longer recognizes,
and not an orphan either), which is why the fidelity section verifies I2
through I1's own check rather than a separate one.

\paragraph{Invariant 3 --- no two hub winners.}
\[
  \# hubOwner \leq 1.
\]
\texttt{HubIdentify} always sets \texttt{hubOwner'} to the singleton
\texttt{\{c?\}}, replacing whatever was there rather than adding to it; no
other operation adds to \texttt{hubOwner}.  So \texttt{hubOwner} can never
hold more than the most recent identify's connection.

\paragraph{Invariant 4 --- deadlock freedom.}
Every operation is enabled for some argument in every reachable state:
\texttt{HubIdentify} and \texttt{ConnectionDrop} carry no precondition
beyond their connection argument's type, so at least one is always
enabled.  No interleaving can trap the model partway through processing a
preemption --- eviction from \texttt{hubOwner} and \texttt{liveConns} and
the new owner's registration happen in the same atomic
\texttt{HubIdentify} step, so there is no intermediate state in which an
old and new claimant are both live.

% ============================================================
\section{Verification}
\label{sec:verify}
% ============================================================

Type-checking is a \texttt{fuzz} pass:
\begin{verbatim}
  fuzz docs/hub_display_reconciliation.tex
\end{verbatim}

I1, I2, and I3 are state properties, checked by asking ProB whether their
negation is \emph{reachable}.  A goal reported \emph{not found} means the
invariant holds on every reachable state.  All three were run and returned
\emph{goal not found} against \texttt{DEFAULT\_SETSIZE 3}:
\begin{verbatim}
  # Invariant 1 (goal NOT found)
  probcli docs/hub_display_reconciliation.tex -model_check -nodead \
    -goal "manifestDeclared = set &
           (card({s | s : SCENE_ID & s : dom(sceneOwner) &
                 not(sceneOwner(s) = orphan) &
                 not(card(hubOwner /\ {c | c : CONNECTION_ID &
                       sceneOwner(s) = conn(c)}) > 0) }) > 0
           or
           card({s | s : SCENE_ID & s : dom(sceneOwner) &
                 sceneOwner(s) = orphan & not(s : manifested)}) > 0)" \
    -p DEFAULT_SETSIZE 3

  # Invariant 2's own direct check (goal NOT found -- see the note under
  # Invariant 2 in Section 6: this is structurally tautological in both
  # models, since ConnectionDrop orphans exactly what it drops liveness
  # for; it is run here as a sanity check on the algebra, not as the
  # falsifiable test -- Invariant 1's check above is what actually
  # distinguishes a correct HubIdentify from the fidelity variant)
  probcli docs/hub_display_reconciliation.tex -model_check -nodead \
    -goal "card({s | s : SCENE_ID & s : dom(sceneOwner) &
                 card({c | c : CONNECTION_ID & sceneOwner(s) = conn(c) &
                       not(c : liveConns)}) > 0}) > 0" \
    -p DEFAULT_SETSIZE 3

  # Invariant 3 (goal NOT found)
  probcli docs/hub_display_reconciliation.tex -model_check -nodead \
    -goal "card(hubOwner) > 1" \
    -p DEFAULT_SETSIZE 3
\end{verbatim}

I4 is the full-state-space deadlock check.  Every state has at least one
enabled operation (Section~\ref{sec:inv}), so any deadlock reported would
be a genuine defect in the model, not a deliberate terminal state.  Run
against \texttt{DEFAULT\_SETSIZE 3}, ProB visits every reachable state and
reports no deadlock:
\begin{verbatim}
  # Invariant 4 / deadlock freedom (no deadlock, all states visited)
  probcli docs/hub_display_reconciliation.tex -model_check \
    -p DEFAULT_SETSIZE 3
\end{verbatim}

The reachable-behaviour witnesses --- preemption actually orphaning a
predecessor's scenes, a manifest actually changing which scenes the
Display holds, and an orphan actually surviving a manifest that still
names it --- must each be \emph{found}, or the invariants above would
hold only because the mechanism they protect is unreachable.  All three
were run and returned \emph{goal found}:
\begin{verbatim}
  # Preemption orphans a predecessor's scene (goal FOUND)
  probcli docs/hub_display_reconciliation.tex -model_check -nodead \
    -goal "card(hubOwner) = 1 &
           card({s | s : SCENE_ID & s : dom(sceneOwner) &
                 sceneOwner(s) = orphan & not(s : manifested)}) = 0 &
           card({s | s : SCENE_ID & s : dom(sceneOwner) &
                 sceneOwner(s) = orphan}) > 0 &
           manifestDeclared = clear" \
    -p DEFAULT_SETSIZE 3

  # A manifest establishes a single orphaned, still-manifested survivor
  probcli docs/hub_display_reconciliation.tex -model_check -nodead \
    -goal "card(dom(sceneOwner)) = 1 &
           card({s | s : SCENE_ID & s : dom(sceneOwner) &
                 sceneOwner(s) = orphan}) = 1 &
           manifestDeclared = set & card(manifested) = 1" \
    -p DEFAULT_SETSIZE 3

  # An orphan survives because the manifest still names it
  probcli docs/hub_display_reconciliation.tex -model_check -nodead \
    -goal "card({s | s : SCENE_ID & s : dom(sceneOwner) &
                 sceneOwner(s) = orphan & s : manifested}) > 0" \
    -p DEFAULT_SETSIZE 3
\end{verbatim}

% ============================================================
\section{Fidelity: reproducing the shipped defect}
\label{sec:fidelity}
% ============================================================

A model that cannot reproduce the known defect proves nothing.  The buggy
variant, kept as
\texttt{docs/hub\_display\_reconciliation\_no\_preemption\_buggy.tex},
removes only the socket-eviction half of preemption from
\texttt{HubIdentify}: \texttt{hubOwner} still correctly replaces (so I3
stays intact --- this variant isolates I1/I2 from I3, rather than
reproducing all three failures at once), but the predecessor's socket is
never closed and its scenes are never orphaned:

\begin{verbatim}
  % correct (HubIdentify in hub_display_reconciliation.tex):
  liveConns' = (liveConns \ hubOwner) \cup {c?}
  sceneOwner' = sceneOwner \oplus (staleScenes \cross {orphan})

  % buggy (this variant):
  liveConns' = liveConns \cup {c?}   % old socket never closes
  sceneOwner' = sceneOwner           % no orphaning at all
\end{verbatim}

This is precisely the interleaving hazard Ground~1 of the design's z-spec
section names: the old connection's socket --- and hence its ability to
still deliver a buffered message --- outlives the identify that was
supposed to have superseded it.

\paragraph{The trace.}
\begin{enumerate}
  \item \texttt{HubIdentify}(old): \texttt{liveConns = \{old\}},
    \texttt{hubOwner = \{old\}}.
  \item \texttt{SceneInstall}(old, s1): guarded on \texttt{old \in
    liveConns}, which holds --- \texttt{sceneOwner(s1) = conn(old)}.
  \item \texttt{HubIdentify}(new) --- buggy: \texttt{hubOwner = \{new\}}
    (still correctly replaced), but \texttt{liveConns = \{old, new\}} ---
    \texttt{old} is never evicted --- and \texttt{s1} is left untouched
    (no orphaning ran).
  \item \texttt{HubManifest}(new, \{\}) --- \texttt{new} declares nothing
    live.  Correctly guarded (\texttt{new \in hubOwner}), so it fires.
    The purge keeps \texttt{s1} only if \texttt{sceneOwner(s1) =
    conn(new)}, which it is not, and \texttt{s1} is not in the empty
    manifest --- so \texttt{s1} is purged.  \texttt{sceneOwner}'s domain
    no longer contains \texttt{s1}.  \texttt{manifestDeclared = set}.
  \item \texttt{SceneInstall}(old, s1) --- the straggler, buffered from
    before step 4, is processed now.  Its guard is \texttt{old \in
    liveConns}: under the buggy variant \texttt{old} is still there, so
    the guard does not disable it --- \texttt{sceneOwner(s1) = conn(old)}
    once more.  The ghost the manifest just purged is back.
\end{enumerate}

At the end of this trace, \texttt{hubOwner = \{new\}},
\texttt{manifestDeclared = set}, and \texttt{sceneOwner(s1) = conn(old)}:
neither owned by the live \texttt{hubOwner} (\texttt{old \neq new}) nor an
orphan, so I1's negation is satisfied.  Step 5 is also the I2 violation
directly --- an install accepted from a connection a correct preemption
would already have evicted --- but I2's own direct reachability check
(``\texttt{sceneOwner} names a connection outside \texttt{liveConns}'')
does \emph{not} distinguish the two models (Section~\ref{sec:inv}'s note
on Invariant~2): \texttt{old} is still in \texttt{liveConns} in the buggy
variant precisely \emph{because} it was never evicted, so that check
finds nothing to report in either spec.  I1's check is what carries the
falsifiable teeth for this defect; running it is how both invariants are
verified against this variant:
\begin{verbatim}
  # must be NOT found for the correct spec, FOUND for the buggy variant
  probcli <spec>.tex -model_check -nodead \
    -goal "manifestDeclared = set &
           (card({s | s : SCENE_ID & s : dom(sceneOwner) &
                 not(sceneOwner(s) = orphan) &
                 not(card(hubOwner /\ {c | c : CONNECTION_ID &
                       sceneOwner(s) = conn(c)}) > 0) }) > 0
           or
           card({s | s : SCENE_ID & s : dom(sceneOwner) &
                 sceneOwner(s) = orphan & not(s : manifested)}) > 0)" \
    -p DEFAULT_SETSIZE 3
\end{verbatim}
With single-owner preemption intact, step 3 evicts \texttt{old} from
\texttt{liveConns} in the same step it replaces \texttt{hubOwner}; by step
5 \texttt{old} can never fire \texttt{SceneInstall} again (I2's own
guard), so the straggler cannot land, and I1's goal returns \emph{not
found} against the correct spec.  Against the buggy variant it returns
\emph{found} --- the trace above.  That difference is the evidence that
the model detects exactly the bug single-owner preemption exists to
close: not merely that two connections can never both claim the identity
(I3, which this variant leaves untouched, isolating the finding), but
that a preempted connection's straggling message can silently
re-materialize a ghost the manifest just told the Display to forget.

% ============================================================
\section{Test Partitions}
\label{sec:partitions}
% ============================================================

Derived from the operations and invariants above, audited against
\texttt{tests/display/test\_hub\_reconciliation.py},
\texttt{tests/display/test\_render\_loop.py},
\texttt{tests/display/replica/test\_scene\_replica.py}, and
\texttt{tests/integration/test\_hub\_display\_reconciliation.py}:

\begin{description}
  \item[P1 --- preemption evicts a live predecessor.] A second
    \texttt{kind="hub"} identify forcibly disconnects the first before
    recording itself.  Covered:
    \texttt{test\_a\_second\_hub\_identify\_forcibly\_disconnects\_the\_first}.
  \item[P2 --- preemption is a no-op with no predecessor.] The ordinary
    restart case --- the old connection's socket is already gone.
    Covered: \texttt{test\_a\_hub\_identify\_with\_no\_predecessor\_preempts\_nothing}.
  \item[P3 --- a \texttt{kind="test"} identify never preempts.] Covered:
    \texttt{test\_a\_test\_identify\_never\_preempts\_or\_marks\_hub},
    \texttt{test\_a\_different\_named\_hub\_identify\_is\_not\_preempted}.
  \item[P4 --- a scene outside the manifest and not fd-owned is purged.]
    Covered: \texttt{test\_a\_scene\_outside\_the\_manifest\_is\_purged},
    \texttt{TestScenesToPurge.test\_a\_scene\_outside\_the\_manifest\_and\_owner\_is\_a\_candidate}.
  \item[P5 --- a manifested scene survives.] Covered:
    \texttt{test\_a\_scene\_in\_the\_manifest\_survives}.
  \item[P6 --- an fd-owned scene survives regardless of the manifest.]
    Covered: \texttt{test\_a\_scene\_owned\_by\_the\_identifying\_fd\_survives}.
  \item[P7 --- a mixed frame loses only its ghost scene, per scene not
    per frame.] Covered:
    \texttt{test\_a\_mixed\_frame\_only\_loses\_its\_ghost\_scene},
    \texttt{TestScenesToPurge.test\_a\_mixed\_frame\_loses\_only\_its\_ghost\_scene}.
  \item[P8 --- an orphan is swept by the same rule as any other ghost, no
    special case.] Covered:
    \texttt{TestScenesToPurge.test\_an\_orphaned\_scene\_is\_swept\_by\_the\_same\_rule}.
  \item[P9 --- a purge closes an emptied frame, silently.] Covered:
    \texttt{test\_a\_manifest\_driven\_purge\_sends\_no\_frame\_close\_event}.
  \item[P10 --- a user-initiated close still notifies (regression guard).]
    Covered: \texttt{test\_a\_user\_initiated\_close\_still\_notifies\_the\_owner}.
  \item[P11 --- end-to-end restart: a second Hub-identified client with an
    empty manifest purges the first connection's scene over the real
    socket.] Covered:
    \texttt{tests/integration/test\_hub\_display\_reconciliation.py}.
\end{description}

\end{document}
